% 乘积法则（综述）
% license CCBYSA3
% type Wiki

本文根据 CC-BY-SA 协议转载翻译自维基百科\href{https://en.wikipedia.org/wiki/Product_rule}{相关文章}

在微积分中，乘积法则（或称莱布尼茨法则[2]、莱布尼茨乘积法则）是一个用于求两个或多个函数乘积的导数的公式。对于两个函数，可以用拉格朗日记号表示为：
$$
(u \cdot v)' = u' \cdot v + u \cdot v',~
$$
或用莱布尼茨记号表示为：
$$
\frac{d}{dx}(u \cdot v) = \frac{du}{dx} \cdot v + u \cdot \frac{dv}{dx}.~
$$
该法则可以推广到三个或更多函数的乘积，推广到乘积的高阶导数规则，以及其他语境中。
\subsection{发现}
这一法则的发现通常归功于 戈特弗里德·莱布尼茨，他使用“无穷小”（现代微分的前身）进行了论证。[3]（然而，莱布尼茨论文的译者 J. M. Child[4] 认为应归功于艾萨克·巴罗。）以下是莱布尼茨的推理：[5]设$u$和$v$是函数。则$d(uv)$等于两个相继$uv$之差；其中一个是$uv$，另一个是$(u+du)(v+dv)$。于是：
$$
d(u \cdot v) = (u+du)\cdot(v+dv) - u\cdot v 
= u \cdot dv + v \cdot du + du \cdot dv.~
$$
由于项$du \cdot dv$相较于$du$和$dv$是“可忽略的”，莱布尼茨得出结论：
$$
d(u \cdot v) = v \cdot du + u \cdot dv,~
$$
这确实就是乘积法则的微分形式。

如果两边同时除以微分$dx$，就得到：
$$
\frac{d}{dx}(u \cdot v) = v \cdot \frac{du}{dx} + u \cdot \frac{dv}{dx},~
$$
这同样可以用拉格朗日记号写作：
$$
(u \cdot v)' = v \cdot u' + u \cdot v'.~
$$
